\subsection{When do the modifications help?}
\owner{Member 3; this paragraph integrates the available Member 3 evidence.}
Tables~\ref{tab:m3-main} and~\ref{tab:m3-diag} distinguish a mechanism benefit from a
practical win. Using a previous second moment reduces the bias observed in the asymmetric-noise
diagnostic at $\beta_2=.9$, including against a control with matched clipping. At the main
setting $\beta_2=.999$, the diagnostic gap is small and ADOPT-0 does not improve held-out
quality on any tested task. Initialization, clipping and transient dynamics provide alternative
explanations for the short-training results; the neural experiments do not isolate these factors.
There is no reduction in tuning effort: both methods receive the same three trials per task.

These conclusions apply to the tested momentum-free specialization, small models, narrow
learning-rate grid and short budget. Image subsets are small; text evaluation covers overlapping
prefix windows; only three main seeds are used. The full paper's default ADOPT, larger models,
longer horizons and complete held-out text regions could behave differently. Results from
Members 1, 2 and 4 must be added before drawing a conclusion across all eight optimizers;
Member 4 supplied configuration and shared code, not measured results for that comparison.
